% Original analytic teaching diagram, 2026-11-10.
\definecolor{VB}{HTML}{7188A3}\definecolor{VBF}{HTML}{E5EAF0}
\definecolor{VG}{HTML}{829C8B}\definecolor{VGF}{HTML}{E8EEE9}
\definecolor{VY}{HTML}{B4A16F}\definecolor{VYF}{HTML}{F1ECDD}
\definecolor{VR}{HTML}{A7787D}\definecolor{VRF}{HTML}{F0E5E6}
\definecolor{VP}{HTML}{9588A5}\definecolor{VPF}{HTML}{ECE8F0}
\definecolor{VI}{HTML}{3B4252}\definecolor{VM}{HTML}{6B7280}\definecolor{VL}{HTML}{4C4D4F}
\begin{tikzpicture}[x=1mm,y=1mm,>=Stealth,font=\figurefont\fontsize{8}{10}\selectfont,text=VI,
lab/.style={text=VM,font=\figurefont\fontsize{7.5}{9}\selectfont}]
% Parabolic boundary of a space-time rectangle; the top is excluded.
\fill[VBF] (0,0) rectangle (60,32);
\draw[VR,line width=2pt] (0,32) -- (0,0) -- (60,0) -- (60,32);
\draw[VM,line width=.9pt,dashed] (0,32) -- (60,32);
\node[below] at (30,-1) {初值：$t=0$};
\node[anchor=east] at (-2,16) {边值};
\node[anchor=west] at (62,16) {边值};
\node[above] at (30,32.6) {$t=T$：不属于抛物边界};
\fill[VI] (30,32) circle (.9);
\draw[VL,->,line width=.8pt] (40,13) -- (31,30.5);
\node[lab,anchor=east] at (57,10) {顶面的极大点也满足 $\partial_t v\geq 0$};
\node at (14,20) {$\Omega\times(0,T]$};
\draw[VL,->,line width=.75pt] (-16,-1) -- (-16,36) node[above] {$t$};
\end{tikzpicture}
